%% ============================================================
%%  ملخص (arabe)
%%  Nécessite XeLaTeX + polyglossia + police arabe (Arial, Amiri...)
%% ============================================================
\chapter*{\texorpdfstring{\textarabic{ملخص}}{Molakhas}}
\addcontentsline{toc}{chapter}{\texorpdfstring{\protect\textarabic{ملخص}}{Molakhas}}
\thispagestyle{plain}

\begin{Arabic}
\begin{flushright}
\setlength{\parindent}{0pt}

يعاني تداول الأوراق المالية الكمي في بورصة الدار البيضاء (مؤشر MASI) من تحديين هيكليين:
تحديد أفضل الفرص السوقية بسرعة في بيئة غنية بالمعلومات، والتحقق الصارم من استراتيجيات
التداول بعيداً عن ظاهرة الإفراط في التوافق (Overfitting).
يهدف هذا المشروع، المنجز ضمن قسم تداول الأسهم في \textbf{BMCE Capital}،
إلى تصميم وتطوير منصة كمية متكاملة منظمة على شكل أربع محطات إنتاج.

\medskip

تتضمن المنصة محركَي نتائج متكاملَين:
الأول، محرك الإشارات (\textit{Signal Engine})، يُقيّم 20 عائلة من المؤشرات التقنية
على 93 ورقة مالية عبر أنبوب معالجة من سبع طبقات (A إلى G)،
مبني على التقييم خارج العينة (OOS) بنوافذ متحركة، ومعيار متعدد المكونات للمتانة،
وتقليص التكرار بالتجميع الارتباطي.
الثاني، محرك الاختبار الراجع (WFO)، يُطبق منهجية التحليل المتقدم لـ Pardo (2008)
مع دالة الهدف PROM، لإنتاج حكم موضوعي على متانة الاستراتيجية عبر معيار Walk-Forward Efficiency.

\medskip

يُغذّي كلا المحركين \textbf{لوحة قيادة موحدة} تمكّن المتداولين من الانتقال
من قراءة السوق إلى اتخاذ قرار الاستثمار في أقل من 30 ثانية لكل ورقة مالية.

\bigskip
\noindent\textbf{الكلمات المفتاحية:}
الاختبار الراجع الكمي، التحليل المتقدم، محرك الإشارات، المؤشرات التقنية،
التقييم عبر الفترات الفرعية، مؤشر MASI، التداول الخوارزمي.

\end{flushright}
\end{Arabic}
